% !TeX root = ../main.tex
\section{Introduction}\label{sec:intro}

Poker occupies a unique position at the intersection of probability theory, game theory, and artificial intelligence. Unlike chess or Go, it is a game of \emph{imperfect information}: each player sees only a fragment of the true state, and the deliberate concealment of cards creates strategic phenomena---bluffing, trapping, and the delicate balance between exploiting weak opponents and remaining unexploitable oneself---that have no analogue in perfect-information games. It is also a game of chance at the level of a single hand yet a game of skill in the long run, a duality that makes it an ideal laboratory for studying decision making under uncertainty. These properties led the artificial intelligence community to adopt Texas Hold'em, the most popular variant, as a benchmark of progressively increasing difficulty: first the limit (fixed-bet) heads-up form, then no-limit heads-up, and finally no-limit multiplayer games against several simultaneous opponents.

The scientific arc of this program is remarkably complete. Early systems of the 1990s and 2000s relied on hand-crafted expert knowledge and opponent modeling~\cite{billings2002challenge}. The arrival of \emph{counterfactual regret minimization} (CFR) in 2007 provided the first practical algorithm for approximating Nash equilibria in enormous imperfect-information games~\cite{zinkevich2007regret}, and by 2015 the Cepheus project had essentially solved heads-up limit Hold'em~\cite{bowling2015ceph}. In 2017 two independent systems, DeepStack~\cite{moravcik2017deepstack} and Libratus~\cite{brown2017libratus}, defeated human professionals in heads-up no-limit Hold'em---the game tree of which contains roughly $10^{160}$ decision points~\cite{sandholm2015abstraction}. Two years later Pluribus extended superhuman play to six-player no-limit Hold'em~\cite{brown2019pluribus}, a setting where the very notion of a unique optimal strategy dissolves and agents must instead target robust, unexploitable policies.

This paper is written for readers who want both the mathematics and the algorithms in one place, at full rigor but without unexplained leaps. Three goals organise the presentation. First, we build the probabilistic core that any strong player or program must internalize: hand combinatorics (Section~\ref{sec:probability}), equity and pot odds, expected value, and the variance that governs bankroll requirements. Second, we develop the algorithmic ladder in order of conceptual depth: rule-based expert systems and Monte Carlo estimation (Section~\ref{sec:evaluation}), the extensive-form game-theoretic framework (Section~\ref{sec:gametheory}), regret minimization and CFR with its modern variants (Section~\ref{sec:cfr}), and reinforcement learning in imperfect-information settings, including the architectures of DeepStack, Libratus, and Pluribus (Section~\ref{sec:rl}). Third, we ground the entire discussion in a working implementation: a strictly rule-compliant No-Limit Hold'em engine with five AI paradigms at one table, which we use in Section~\ref{sec:experiments} to quantify---via a 100{,}000-hand simulation---how abstract theory translates into observable differences in win rate and playing style.

The paper is deliberately self-contained. Every algorithm described is either derived from first principles or accompanied by pseudocode faithful to a working implementation; every probabilistic claim is either proved or reducible to a one-line application of a stated theorem. Where a rigorous treatment would require measure-theoretic machinery---for instance, the full generality of stochastic games---we restrict attention to the finite structures that Hold'em actually presents, which suffice for every result we state. Readers interested primarily in the practical study can jump directly to Section~\ref{sec:experiments} after Sections~\ref{sec:gametheory} and~\ref{sec:cfr}, on which the experimental agents' designs rest.
